\newpage
\section{UVX Guide}

Using Python's toolchain runner (\textbf{uvx}) alongside Node gives you access to a completely different class of system automation tools. Running servers via \lstinline[style=inline]|uvx| provides lightning-fast local execution directly from registries.

\subsection{Official File System and Git Servers}
These servers give the agent native capabilities on the local machine without spawning blind bash wrappers.

\begin{itemize}
    \item \textbf{Git History} (\lstinline[style=inline]|mcp-server-git|): Gives the agent native capabilities to read branch histories, calculate diffs, stage files, and check commits.
    \item \textbf{Local Filesystem} (\lstinline[style=inline]|mcp-server-filesystem|): Opens up file I/O operations (reading directories, verifying structures, changing source code) across any system root you expose to it.
    \item \textbf{Fetch Markdown} (\lstinline[style=inline]|mcp-server-fetch|): A Python alternative to the npx fetch server.
\end{itemize}

\subsection{Local Custom Servers with Editable Installs}
Unlike Node, where \lstinline[style=inline]|npm link| is required for seamless path resolution, \textbf{uv} natively understands local editable development and workspaces. You can tell \lstinline[style=inline]|uvx| to run your local packages directly via their file paths without needing a separate global linking step.

\paragraph{Setting up the uv Workspace.}
Your root \lstinline[style=inline]|pyproject.toml| should explicitly declare your subdirectories as part of a \textbf{uv} workspace:

\begin{lstlisting}[style=toml]
[project]
name = "my-mcp-suite"
version = "0.1.0"
requires-python = ">=3.10"

[tool.uv.workspace]
members = ["packages/*"]
\end{lstlisting}

Inside your server subdirectory's \lstinline[style=inline]|pyproject.toml|, ensure you expose the script entry point:

\begin{lstlisting}[style=toml]
[project.scripts]
my-filesystem-package-name = "service_filesystem.main:run"
\end{lstlisting}

\paragraph{Running Editable Installs.}
Configure your agent to run the server using the \lstinline[style=inline]|--with editable+| flag. When the agent triggers the server, \lstinline[style=inline]|uv| will instantly run the live code from your directory, applying code changes instantly without rebuilds.

\begin{lstlisting}[style=json]
{
  "mcpServers": {
    "my-local-filesystem": {
      "command": "uvx",
      "args": [
        "--with",
        "editable+/absolute/path/to/my-mcp-suite/packages/service-filesystem",
        "my-filesystem-package-name"
      ]
    }
  }
}
\end{lstlisting}

\subsection{Universal UVX Configuration Block}
When configuring these in your runtime hub (e.g., \lstinline[style=inline]|settings.json| or \lstinline[style=inline]|mcp_config.json|), you can seamlessly combine your local editable Python packages alongside official distributions.

\begin{lstlisting}[style=json]
{
  "mcpServers": {
    "git-history": {
      "command": "uvx",
      "args": ["mcp-server-git"]
    },
    "local-filesystem": {
      "command": "uvx",
      "args": ["mcp-server-filesystem", "/"]
    },
    "my-local-filesystem": {
      "command": "uvx",
      "args": [
        "--with",
        "editable+$WORKSPACE_DIR/packages/service-filesystem",
        "my-filesystem-package-name"
      ]
    }
  }
}
\end{lstlisting}
